\section{Predikat dan Arity (Relasi $n$-ary)}

Jika term berperan sebagai argumen yang dibicarakan, predikat menyatakan sifat atau relasi yang dikenakan pada term tersebut. Tanpa predikat, rangkaian term tidak membentuk klaim yang dapat dievaluasi. Predikat-lah yang, bersama interpretasi semesta, menentukan apakah suatu kalimat tertutup \textbf{benar (T)} atau \textbf{salah (F)} \cite{rosen2019}.

\subsection{Definisi Formal Predikat}

\begin{definisi}[Predikat]
\logic{Predikat} merujuk pada atribut, klasifikasi, properti, kondisi, atau relasi logis yang dilekatkan pada satu atau lebih term.
\end{definisi}

Menurut konvensi, predikat sering dinotasikan dengan huruf kapital, misalnya $P$, $Q$, $R$, $S$, atau simbol deskriptif seperti $M$ untuk ``manusia''.

\begin{contoh}
Contoh klasik: ``Socrates adalah manusia''.
\begin{enumerate}
    \item Term (konstanta): $s \equiv \text{Socrates}$
    \item Bagian predikatif: ``adalah seorang manusia''
    \item Predikat: $M(x) \equiv x \text{ merupakan manusia}$.
\end{enumerate}

Gabungan subjek--predikat dalam notasi formal: $M(s)$. Jika dalam interpretasi yang dimaksud Socrates adalah manusia, maka $M(s)$ bernilai \textbf{benar (T)}.
\end{contoh}

\subsection{Arity (banyaknya argumen)}

Banyak sifat hanya menyangkut satu objek; relasi lain memerlukan dua objek atau lebih. \textbf{Arity} predikat adalah banyaknya posisi argumen (tempat term) yang disediakan oleh simbol predikat tersebut.

\subsubsection{Predikat 1-arity (\textit{Unary} / monadik)}
Mengikat satu term; lazim untuk properti tunggal.
\begin{itemize}
    \item $G(n) \equiv n \text{ adalah bilangan genap}$
    \item $W(x) \equiv x \text{ merupakan warga benua Afrika}$
    \item Contoh evaluasi: $G(8)$ bernilai \textbf{True}; $G(13)$ bernilai \textbf{False}.
\end{itemize}

\subsubsection{Predikat 2-arity (\textit{Binary} / diadik)}
Mengikat sepasang term. Urutan argumen penting pada relasi tidak simetris.
\begin{itemize}
    \item $L(x, y) \equiv x \text{ lebih besar daripada } y$.
    \item $S(a, b) \equiv a \text{ memiliki hubungan asmara terhadap } b$ (contoh relasi naratif).
    \item $L(100, 2)$ bernilai \textbf{True}, sedangkan $L(2, 100)$ bernilai \textbf{False}.
    \item Untuk relasi seperti $S$, dari $S(\text{Romeo}, \text{Juliet})$ belum tentu diperoleh $S(\text{Juliet}, \text{Romeo})$; asimetri relasi tercermin dalam urutan argumen \cite{wiki_first_order_logic}.
\end{itemize}

\subsubsection{Predikat $n$-arity (\textit{n-ary} / poliadik)}
Predikat dapat memiliki lebih dari dua argumen.
\begin{itemize}
    \item \textbf{Ternary}: $B(x, y, z)$ dapat bermakna ``kota $x$ terletak pada garis lurus antara kota $y$ dan kota $z$''.
    \item Contoh ilustratif: $B(\text{Salatiga}, \text{Semarang}, \text{Boyolali})$ dapat dianggap benar dalam model geografis sederhana yang disepakati.
\end{itemize}
